\section{Introduction}
\label{sec:intro}

Selective visual attention requires a behaving animal to perform two logically distinct computations concurrently. It must form a \emph{policy} --- a moment-by-moment decision about what to do, such as whether to report that a behaviourally relevant event has occurred --- and it must form a \emph{value estimate}, a prediction of how good its current situation is, which under any reward-driven account is precisely the quantity that supplies the teaching signal for the policy. These two computations draw on overlapping but distinguishable information. \emph{When and where to act} is driven by bottom-up, image-grounded sensory evidence accumulated toward a commitment threshold \citep{gold2007neural}; \emph{how valuable the moment is} depends on accumulated context --- which location was cued, how reliable the cue is, how much reward is at stake --- a quantity carried on the same parietal substrate that supports the decision itself \citep{platt1999neural, sugrue2004matching}. This decomposition is the actor--critic organisation of reinforcement learning, and it has a deep counterpart in the biology: an actor for the policy and a critic for state value are carried by anatomically distinct but functionally interdependent circuits \citep{odoherty2004dissociable, maia2009reinforcement}.

\paragraph{The design question.}
Given an attention agent that must compute both a policy and a value, how should the two computations be wired? A natural engineering instinct is to \emph{separate} them: give each its own specialised recurrent stream so that the image-reading machinery that decides when to act is insulated from the context-integrating machinery that estimates value. Separation is attractive on grounds of modularity and interpretability, and for purely feedforward representations there is evidence that it helps. The alternative is to \emph{couple} the streams through a shared, jointly maintained working memory, so that the actor writes the very state the critic evaluates. Which organisation is required --- and whether the question even has a categorical answer --- is the subject of this paper.

\paragraph{A lesion-style controlled comparison.}
We train a conv-free recurrent vision-transformer agent purely from sparse reward on a spatially cued orientation-change-detection task ($T=29$ frames) built to mirror primate covert-attention paradigms. The agent embeds each $50\times50$ frame as $100$ patch tokens on a $10\times10$ grid, maintains per-token recurrent memories through a cross-attention encoder, and reads out two heads: an actor that emits a binary wait/declare action each frame and a distributional (quantile) critic. Two recurrent streams feed these heads --- a \emph{priority stream} ($\mu$) whose residual preserves the raw image and which therefore carries live visual evidence to the actor, and a \emph{value stream} ($q$) whose residual is the deep memory and which carries accumulated context to the critic. (In an older vocabulary these are the ``salience'' and ``top-down'' streams; we use priority and value throughout.) We then hold the task, the learning algorithm, the hyperparameters and the number of updates fixed and vary a single structural property across three configurations, in the manner of a lesion experiment. The first \emph{couples} the streams through a shared recurrent memory with the priority stream commanding the actor. The second \emph{severs} the coupling --- each stream is given a private memory updated only by its own output, the minimal edit that removes cross-talk while preserving the split, the stream types, the parameter count and the recipe. The third \emph{retains} the coupling but \emph{swaps the routing}, so the value stream drives the actor and the priority stream drives the critic. Only one wiring can show what is structurally required, because the other two isolate the two candidate necessary conditions.

\paragraph{Preview of the result.}
The outcome is categorical rather than graded. The coupled, correctly routed agent becomes a calibrated psychophysical detector (accuracy $0.868$, hit rate $0.794$, false-alarm rate $0.067$, reaction time $\sim$$1.95$ frames) that reproduces the reliability-scaled validity benefit of primate spatial attention. Severing the coupling collapses the agent to always-wait: across $500$ episodes it never once declares a change (hit rate $0.000$; accuracy $0.504$, exactly the no-change base rate). Swapping the routing collapses it identically (hit rate $0.000$; accuracy $0.420$). Learning therefore requires \emph{both} that the loop be closed through a shared recurrent memory \emph{and} that the grounded priority stream command the policy; a single structural edit to either condition produces a $\sim$$79$-point hit-rate gap. We trace the collapses to reinforcement-learning credit assignment --- a critic that evaluates a state the actor does not occupy cannot teach it, and an actor fed a representation that cannot see the change has nothing to which credit can be assigned.

\paragraph{Brain-loop interpretation and novelty.}
We relate the one working configuration to the primate cortico-basal-ganglia--collicular loop: a grounded priority and commitment node (visual cortex, FEF-visual cells, superior colliculus) that commits the action and writes the shared working state \citep{thompson2005neuronal, moore2001control, lovejoy2010inactivation, zenon2012attention}; a value/critic system in the basal ganglia that reads that state to evaluate it \citep{odoherty2004dissociable, maia2009reinforcement}; and a dopaminergic reward-prediction-error signal from substantia nigra pars compacta that closes the loop by teaching the actor's plasticity \citep{schultz1997neural, hollerman1998dopamine, morita2014striatal}. We present this correspondence as conceptual motivation, not literal isomorphism, and state its limits plainly in the Discussion. The result is, to our knowledge, novel for a recurrent, image-grounded agent: no prior model splits a recurrent vision encoder into priority and value streams on a shared memory, requires that coupling for learning, or requires the grounded stream to be the actor. It stands in clean contrast to the one machine-learning finding on actor/critic representation sharing, which reports the \emph{opposite} for feedforward representations --- there, separation helps. For a recurrent agent that must read a live image to act, coupling is required.
